\chapter{Alias and Pointer Analysis}
\label{ch:ch20-alias-analysis}

\section{Two Names, One Cell}
\label{sec:ch20-two-names}

Every optimization that touches memory must answer one question first.
Do these two references access overlapping bytes?
If \code{*p} and \code{*q} name the same cell, then a store through \code{q} changes the value a later load through \code{p} returns.
If they never overlap, the load is independent of the store and all the machinery of earlier chapters applies unchanged.
Guessing wrong in either direction hurts.
Assuming independence when the locations overlap miscompiles.
Assuming overlap when the locations are disjoint kills motion, numbering, and scheduling for no reason.
\term{Alias analysis} computes, for each pair of references, whether they may, must, or cannot reach the same location.
\term{Pointer analysis} computes, for each pointer, the set of abstract objects it may target.
Intersecting points-to sets indicate possible aliasing; abstraction and lost path correlations can make that intersection a false positive.
Disjoint sound sets prove no alias within the memory model being analyzed.
Precision is undecidable in general, so every analysis approximates.
The rule for soundness is that a \term{may} answer must over-approximate: it must report overlap whenever overlap is possible, even at the cost of false alarms.
A \term{must} answer must under-approximate: it may claim identity only when identity holds on every execution.
This chapter builds a small but honest analysis, classifies its answers, shows how clients consume them, and explains what more expensive analyses buy.
Chapter~\ref{ch:interprocedural-heap-analysis} supplies the semantic background on abstract heap objects and updates; here the emphasis is an inclusion-constraint solver and the alias queries exposed to optimization clients.

\section{May, Must, and Cannot}
\label{sec:ch20-relations}

Distinguish semantic relations from conservative query answers.
Semantically, two accesses may alias when some execution makes their byte ranges overlap.
An analysis answer of \term{may-alias} means that overlap has not been ruled out; it need not supply such an execution.
Here \term{must-alias} means equal starting addresses on every relevant execution; access sizes must be checked separately.
\term{No-alias} holds when no execution overlaps them.
May-alias is the workhorse, because it gates every motion and reuse: a store kills a load exactly when the two may alias.
Must-alias supports forwarding when an ordinary dominating store supplies all bytes of the load, with compatible representation and no intervening clobber.
No-alias is the prize, because it restores the full power of value-based reasoning to memory.
Consider three tiny situations.
After \code{p = \&x}, the expressions \code{*p} and \code{x} must alias, since every execution binds \code{p} to \code{x}.
After a join where one branch sets \code{p = \&x} and the other sets \code{p = \&y}, the pair \code{*p} and \code{x} may alias but does not must alias.
After \code{p = \&x} followed by \code{q = \&y} for distinct locals \code{x} and \code{y}, the pair \code{*p} and \code{*q} cannot alias.
Distinct locals, distinct globals, and distinct heap allocation sites give the analysis its base facts of separation.
Everything else is propagation of those facts through assignments, joins, and calls.
Two design axes separate the many published analyses.
\term{Flow sensitivity} distinguishes program points, so a pointer can target \code{x} before an assignment and \code{y} after it.
\term{Context sensitivity} distinguishes call sites, so a helper called with \code{\&x} in one place and \code{\&y} in another keeps the two cases apart.
Flow-insensitive and context-insensitive analyses merge all of this information into one global answer per pointer.
They are cheaper and less precise, and they are the right starting point for a first implementation.
Chapter~\ref{ch:ch21-memory-optimization} builds its load and store optimizations directly on the answers defined here.
% TODO(cite): primary references for inclusion-based and unification-based pointer analyses.

\section{A Flow-Insensitive Inclusion Analysis}
\label{sec:ch20-inclusion}

The simplest useful analysis treats the program as a bag of constraints and solves them to a fixed point.
Each variable \code{v} gets a \term{points-to set} $\mathit{pts}(v)$ holding the abstract objects \code{v} may target.
Each kind of statement contributes one constraint.
An address assignment \code{p = \&x} adds the object \code{x} to $\mathit{pts}(p)$.
A copy \code{p = q} adds everything in $\mathit{pts}(q)$ to $\mathit{pts}(p)$.
A load \code{p = *q} adds, for every object \code{o} in $\mathit{pts}(q)$, everything in $\mathit{pts}(o)$ to $\mathit{pts}(p)$.
A store \code{*p = q} adds, for every object \code{o} in $\mathit{pts}(p)$, everything in $\mathit{pts}(q)$ to $\mathit{pts}(o)$.
The inclusion direction matters: information flows from right to left along assignments, and sets only grow.
Because sets only grow, iterating the constraints in any order reaches the same least fixed point.
Listing~\ref{lst:ch20-constraints} gives the example we will solve by hand.

\begin{irlisting}[caption={Five statements whose points-to constraints interact through two copies.},label={lst:ch20-constraints}]
p = &x
q = &y
p = q
r = p
s = &x
\end{irlisting}

The constraints read directly off the statements.
First, $\{x\} \subseteq \mathit{pts}(p)$, $\{y\} \subseteq \mathit{pts}(q)$, and $\{x\} \subseteq \mathit{pts}(s)$.
Second, $\mathit{pts}(q) \subseteq \mathit{pts}(p)$ from the copy \code{p = q}.
Third, $\mathit{pts}(p) \subseteq \mathit{pts}(r)$ from the copy \code{r = p}.
Algorithm~\ref{alg:ch20-points-to} solves constraints of this form with a worklist.

\begin{algorithm}
  \caption{Worklist solver for inclusion points-to constraints}\label{alg:ch20-points-to}
  \begin{pseudocode}
    \Require finite closed-world constraints over pointer cells and abstract objects
    \Ensure least fixed point of all points-to sets
    \State set every $\mathit{pts}(v)$ to $\emptyset$
    \State place every constraint on the worklist $W$
    \While{$W \neq \emptyset$}
      \State remove a constraint $c$ from $W$
      \State evaluate $c$ using current sets, registering every set read as a dependency
      \State union its contributions into the destination sets
      \If{some target set gained a new object}
        \State add every constraint reading that set to $W$
      \EndIf
    \EndWhile
  \end{pseudocode}
\end{algorithm}

The algorithm assumes a closed world of abstract objects: locals, globals, and one object per allocation site.
Each pointer-valued object also has a points-to set describing its stored contents.
For a load \code{p = *q}, register reads of $\mathit{pts}(q)$ and every $\mathit{pts}(o)$ for $o\in\mathit{pts}(q)$.
If $q$ later gains a target, revisit the load, register the new dependency, and immediately propagate the target's existing contents.
For a store \code{*p = q}, register reads of both pointer sets and update every object currently in $\mathit{pts}(p)$.
Missing a newly discovered dependency can leave a load permanently stale even after its target's contents grow.
Keeping obsolete dependency registrations is safe; it only causes extra work.
Unknown external pointers require an explicit unknown-object policy or conservative initial sets; an empty set must not stand for an arbitrary external address.
It ignores control flow entirely, including the order of statements and the conditions on branches.
That choice is the source of both its speed and its imprecision.
Termination follows because each set grows monotonically and the universe of objects is finite.
Each constraint re-enters the worklist only when an input set grows, so total work stays proportional to the number of propagations.

\section{Worked Example: Solving the Constraints}
\label{sec:ch20-worked-solving}

We solve Listing~\ref{lst:ch20-constraints} by hand, tracking every set after each round.
Initially all sets are empty: $\mathit{pts}(p) = \emptyset$, $\mathit{pts}(q) = \emptyset$, $\mathit{pts}(r) = \emptyset$, and $\mathit{pts}(s) = \emptyset$.
Round one applies the three address constraints.
The statement \code{p = \&x} sets $\mathit{pts}(p) = \{x\}$.
The statement \code{q = \&y} sets $\mathit{pts}(q) = \{y\}$.
The statement \code{s = \&x} sets $\mathit{pts}(s) = \{x\}$.
For this trace, process only address constraints in round one and defer copies to round two; thus $\mathit{pts}(r)$ is still empty.
Round two propagates the copies with the newly grown inputs.
The copy \code{p = q} adds $\mathit{pts}(q) = \{y\}$ to $\mathit{pts}(p)$, giving $\mathit{pts}(p) = \{x, y\}$.
The copy \code{r = p} adds the current $\mathit{pts}(p) = \{x, y\}$ to $\mathit{pts}(r)$, giving $\mathit{pts}(r) = \{x, y\}$.
Round three rechecks every constraint against the final sets.
The copy \code{p = q} contributes nothing new, the copy \code{r = p} contributes nothing new, and the address constraints are already satisfied.
The fixed point is $\mathit{pts}(p) = \{x, y\}$, $\mathit{pts}(q) = \{y\}$, $\mathit{pts}(r) = \{x, y\}$, and $\mathit{pts}(s) = \{x\}$.
Notice the deliberate imprecision: the statement \code{p = q} overwrites \code{p} in execution order, so a flow-sensitive analysis would report $\mathit{pts}(p) = \{y\}$ afterward.
Our flow-insensitive result $\{x, y\}$ keeps the stale \code{x}, which is sound but weaker.
Table~\ref{tab:ch20-queries} turns the fixed point into alias answers for the queries clients actually ask.

\begin{table}
  \centering
  \caption{Alias answers derived from the fixed point by intersecting target sets.}
  \label{tab:ch20-queries}
  \begin{tabular}{lll}
    \toprule
    Query pair & Shared targets & Answer \\
    \midrule
    \code{*p} versus \code{*q} & $y$ & may-alias \\
    \code{*r} versus \code{x} & $x$ & may-alias \\
    \code{*q} versus \code{*s} & none & no-alias \\
    \code{*p} versus \code{*s} & $x$ & may-alias \\
    \bottomrule
  \end{tabular}
\end{table}

The third row shows the analysis earning its keep: \code{q} targets only \code{y} while \code{s} targets only \code{x}, so a store through \code{q} never disturbs a load through \code{s}.
The first row is a genuine possible alias: after \code{p = q}, both point to \code{y}.
The fourth row shows the loss of flow sensitivity: at the listing's end, \code{p} points to \code{y} and \code{s} to \code{x}, but the merged set for \code{p} retains its earlier target.
Figure~\ref{fig:ch20-merge} shows why merging at joins forces this conservatism even in a flow-sensitive setting.

\begin{figure}
  \centering
  \begin{tikzpicture}[node distance=14mm and 20mm]
    \node[basic block] (ch20e) {E: p = \&x; q = \&y};
    \node[basic block,below left=of ch20e] (ch20t) {T: p = q};
    \node[basic block,below right=of ch20e] (ch20f) {F: };
    \node[basic block,below right=of ch20t] (ch20j) {J: r = p};
    \draw[cfg edge] (ch20e) -- (ch20t);
    \draw[cfg edge] (ch20e) -- (ch20f);
    \draw[cfg edge] (ch20t) -- (ch20j);
    \draw[cfg edge] (ch20f) -- (ch20j);
  \end{tikzpicture}
  \caption{The join block sees both bindings of the pointer, so its target set is the union of the two branches.}
  \label{fig:ch20-merge}
\end{figure}

At block $J$, the pointer \code{p} targets \code{y} along the left path and \code{x} along the right path.
Any analysis that merges paths reports $\{x, y\}$, and only a path-sensitive analysis keeps the two cases apart.
Path sensitivity is precise but expensive, so most compilers accept the merge and spend their budget on context sensitivity instead.

\section{Implementation Notes}
\label{sec:ch20-implementation}

Let $P$ count pointer cells, including pointer-valued objects, $O$ count abstract objects, and $C$ count constraints.
At most $PO$ distinct memberships can be added.
A simple reference solver can rescan all constraints until a sweep adds nothing: with $K=\lceil O/w\rceil$ words per set, one sweep costs at most $O(COK)$ for the load and store constraints, and at most $PO+1$ sweeps are needed.
This loose bound belongs to that reference implementation; efficient inclusion solvers use dynamic copy edges and delta propagation rather than repeatedly scanning every target set.
% TODO(cite): published complexity bounds for inclusion-based points-to analysis.
Sparse representations change the constant factors dramatically: bitsets for points-to sets, difference propagation so only new objects travel, and offline constraint-graph collapsing before solving.
Context insensitivity is the main precision lever after sparsity.
Cloning the analysis per call site, or summarizing each procedure once and instantiating the summary per caller, recovers the common case where a helper is called with different objects.
Must-alias and no-alias answers need extra care at scale.
Intersecting may-points-to sets does not establish must-alias: a shared abstract allocation site may represent distinct live objects.
Equal singleton targets justify must-alias only when the target denotes one concrete object and both references select the same offset on every relevant execution.
Type-based separation must follow the source language's access rules, including C's character-type, aggregate, union, and effective-type cases, rather than a blanket test for unequal types.
% TODO(cite): effective-type and aliasing rules in the C specification.
But type rules have exceptions, unsafe casts break them silently, and a wrong no-alias answer miscompiles, so verify the language rules before enabling type-based reasoning.
Heap abstraction is the other scaling decision.
One object per allocation site is the standard choice: precise enough to separate unrelated structures, coarse enough to stay small.
Field sensitivity refines each object by member, which pays off in object-oriented code with many small fields.

\section{Exercises}
\label{sec:ch20-exercises}

\begin{enumerate}
  \item \label{ex:ch20-solve-by-hand}
  Solve the constraints for \code{a = \&m}, \code{b = a}, \code{c = \&n}, \code{*b = c}, and \code{d = *a} with Algorithm~\ref{alg:ch20-points-to}.
  Show every points-to set after each round until the fixed point.
  Give $\mathit{pts}(d)$ and explain which store constraint placed each of its members there.
  \item \label{ex:ch20-must-vs-may}
  Give a three-block CFG where \code{*p} and \code{z} must alias at the end of the left branch but only may alias at the join.
  Compute the points-to sets at each point and at the join, showing the union that weakens the answer.
  State which client optimization the must-alias answer enables on the left branch and why the join answer no longer enables it.
  \item \label{ex:ch20-context-loss}
  Write a procedure \code{set} that stores its pointer argument into a global, called once with \code{\&x} and once with \code{\&y}.
  Run the analysis of this chapter on the whole program and give the points-to set of the global.
  Explain why the answer merges the two call sites, state the alias query this imprecision defeats, and describe what per-caller information would recover it.
\end{enumerate}

\section{Further reading}
\label{sec:ch20-further-reading}

Morgan's \emph{Building an Optimizing Compiler} treats aliasing as a practical constraint on every memory optimization~\cite{morgan-1998-optimizing}.
Read it for how alias answers gate the transformations in the surrounding chapters.
Appel's \emph{Modern Compiler Implementation in ML} presents pointer-related reasoning inside a working compiler framework~\cite{appel-1998-modern}.
Use it to see where the analysis results plug into the rest of the pipeline.
The C11 committee draft defines the object, pointer, and effective-type rules that ground any sound C alias analysis~\cite{wg14-2011-c11}.
Check its clauses on object identity and pointer comparison before trusting a separation fact.
